%(BEGIN_QUESTION)
% Copyright 2010, Tony R. Kuphaldt, released under the Creative Commons Attribution License (v 1.0)
% This means you may do almost anything with this work of mine, so long as you give me proper credit

Et type K-termoelement er satt inn i en prosess, med digitalt multimeter koblet til terminalene. Omgivelsestemperaturen ved DMM-ledningene er 84 $^{o}$F. Beregn måleskjøttemperaturen for følgende millivoltmålinger (avrundet til nærmeste hele grad Fahrenheit):

\begin{itemize}
\item{} 2.55 mV ; $T$ = \underbar{\hskip 50pt} $^{o}$F
\vskip 10pt
\item{} 6.21 mV ; $T$ = \underbar{\hskip 50pt} $^{o}$F
\vskip 10pt
\item{} 10.93 mV ; $T$ = \underbar{\hskip 50pt} $^{o}$F
\vskip 10pt
\item{} 18.83 mV ; $T$ = \underbar{\hskip 50pt} $^{o}$F
\end{itemize}

\underbar{file i00390}
%(END_QUESTION)





%(BEGIN_ANSWER)

Alle svar basert på ITS-90 termoelementtabell:

\begin{itemize}
\item{} 2.55 mV ; $T$ = {\bf 195} $^{o}$F
\vskip 10pt
\item{} 6.21 mV ; $T$ = {\bf 357} $^{o}$F
\vskip 10pt
\item{} 10.93 mV ; $T$ = {\bf 567} $^{o}$F
\vskip 10pt
\item{} 18.83 mV ; $T$ = {\bf 904} $^{o}$F
\end{itemize}

%(END_ANSWER)





%(BEGIN_NOTES)


%INDEX% Measurement, temperature: thermocouple millivoltage interpretation

%(END_NOTES)
